\documentclass[aspectratio=169]{beamer}
\usepackage{tcolorbox}

% Open this file in the TikZ editor's Beamer view.
% Click a grey source card to select its original LaTeX in the source panel.
% Replace the selected environment with plain text, then try Undo.
% On the last two slides, advance the overlay step to reveal the card.

\begin{document}

\begin{frame}[t]{Click a source card}
  This paragraph stays editable on the canvas.

  \begin{tcolorbox}[title=An unsupported environment]
    This box appears as a source card in the editor.
    Click the card to select this entire environment.
  \end{tcolorbox}

  The text after the card stays visible and editable too.
\end{frame}

\begin{frame}[t]{Source cards inside columns}
  \begin{columns}[t]
    \begin{column}{.48\textwidth}
      Text before the card.

      \begin{tcolorbox}[title=Preserved source]
        \begin{itemize}
          \item This nested list belongs to the box.
          \item The whole box is selected together.
        \end{itemize}
      \end{tcolorbox}

      Text after the card.
    \end{column}
    \begin{column}{.48\textwidth}
      This column renders normally.

      \begin{itemize}
        \item Editable list text
        \item Native math: \(E = mc^2\)
      \end{itemize}
    \end{column}
  \end{columns}
\end{frame}

\begin{frame}[t]{Overlay: reserve space}
  Advance to step 2 to reveal the card.

  \begin{uncoverenv}<2>
    \begin{tcolorbox}
      This card is hidden on step 1, but keeps its space.
    \end{tcolorbox}
  \end{uncoverenv}

  This paragraph stays in the same position on both steps.
\end{frame}

\begin{frame}[t]{Overlay: remove and reflow}
  Advance to step 2 to insert the card.

  \begin{onlyenv}<2>
    \begin{tcolorbox}
      This card takes up space only on step 2.
    \end{tcolorbox}
  \end{onlyenv}

  This paragraph moves down when the card appears.
\end{frame}

\end{document}
